% !TEX TS-program = XeLaTeX
% !TEX encoding = UTF-8 Unicode

\chapter*{\hfill 结　　论 \hfill}
\addcontentsline{toc}{chapter}{结　　论}
\label{conclusion}
\defaultfont

本文围绕基于 WiFi CSI 的脊柱姿态检测问题，研究了一种非接触式三分类识别方法。该方法利用人体姿态变化对无线信道的扰动作用，将脊柱姿态检测建模为正常姿态、脊柱后凸相关姿态和脊柱侧弯相关姿态的分类任务，并通过接收天线相位比、文件级统计特征和 ExtraTrees 集成分类器实现姿态识别。

本文首先分析了 WiFi CSI 在人体姿态感知中的作用，说明 CSI 幅度和相位能够反映无线传播路径变化。针对原始相位容易受到硬件偏移和公共相位误差影响的问题，本文采用接收天线相位比构造相对相位特征，并沿时间维度进行相位展开。随后，本文使用滑动窗口组织连续 CSI 序列，并将每个原始文件聚合为 1620 维统计特征。最终，本文采用 ExtraTrees 分类器学习统计特征与脊柱姿态类别之间的映射关系。

实验使用 750 个 CSI 原始文件，三类姿态各 250 个样本，并在文件级划分训练集、验证集和测试集。最终模型在 113 个测试文件上取得 92.04\% 的准确率和 92.08\% 的宏平均 F1 值。其中，正常姿态 F1 值为 92.31\%，脊柱后凸相关姿态 F1 值为 89.47\%，脊柱侧弯相关姿态 F1 值为 94.44\%。实验结果表明，基于相位比统计特征和 ExtraTrees 的方法能够有效区分三类脊柱姿态。

本文建立了从 CSI 数据读取、特征处理、模型训练到检测输出的完整技术流程，为基于 WiFi 信号的脊柱姿态辅助检测提供了一种可行方案。该方法具有非接触、低侵入和隐私友好的特点，适用于室内场景下的人体姿态感知研究。后续工作可从受试者规模扩展、跨人泛化评估、幅度与相位融合、环境增强和实时检测流程设计等方面继续展开，以进一步提升系统在实际应用中的稳定性和适用范围。
